% The prediction ledger's guards and the cluster sign-flip test (build/ledger-guards).
\documentclass[11pt]{article}
\usepackage[T1]{fontenc}
\usepackage{lmodern}
\usepackage{microtype}
\usepackage[margin=1in]{geometry}
\usepackage{amsmath,amssymb,amsthm,mathtools}
\usepackage{algorithm}
\usepackage{algpseudocode}
\usepackage{booktabs}
\usepackage[numbers,sort&compress]{natbib}
\usepackage[hidelinks]{hyperref}
\usepackage[capitalise,noabbrev]{cleveref}

\theoremstyle{plain}
\newtheorem{theorem}{Theorem}[section]
\newtheorem{lemma}[theorem]{Lemma}
\newtheorem{proposition}[theorem]{Proposition}
\newtheorem{corollary}[theorem]{Corollary}
\theoremstyle{definition}
\newtheorem{definition}[theorem]{Definition}
\theoremstyle{remark}
\newtheorem{remark}[theorem]{Remark}

\newcommand{\lean}[1]{\texttt{#1}}
\newcommand{\sgn}{\varepsilon}
\newcommand{\one}{\mathbf{1}}

\title{Guards for a Pre-Registered Prediction Ledger: a Pin Before the First Prediction, One
Reading at $N$, and a Cluster Sign-Flip Test}
\author{mind-palace contributors}
\date{7 November 2026}

\begin{document}
\maketitle

\begin{abstract}
A prediction experiment compares two arms on yes/no questions that resolve later, under a protocol
recorded as the first line of a hash-chained ledger. We close three ways to choose a result after
seeing the data: an arm the primary comparison names must be pinned to a sha256 configuration
digest before its first prediction; once $N$ is fixed, only the first $N$ clusters in ledger order
are scored; and a protocol may declare a cluster sign-flip test, exact for exchangeable arms in
independent clusters, whose inverted interval has weighted subset means as its ends. On simulated
nulls with 40 clusters it rejected 5.05\% at $\alpha = 0.05$, the percentile bootstrap 6.86\%; both
over-reject when clusters are dependent.
\end{abstract}

\section{Introduction}

A prediction experiment, as the ledger described here records it, compares arms (a model, an agent
with and without a tool, one configuration against another) on yes/no questions that resolve
later. Each arm states a probability for each question; the primary comparison is the mean
difference in Brier score between a baseline and an arm on the questions both predicted. Questions
that turn on the same underlying event form a cluster, because their outcomes are not independent,
and the cluster is the unit of inference. The analysis is fixed in advance: the experiment's
ledger is an append-only, hash-chained file whose first record is the protocol, which names the
arms, the comparisons, the smallest effect worth detecting $\delta$, the level $\alpha$, the power,
and the resampling seed. After the first prediction an amendment can add to the protocol but cannot
move that bar.

A bar written in advance still leaves three ways for the result to be chosen after the data are
seen, and this paper closes them:
\begin{enumerate}
\item \emph{A configuration chosen afterwards.} Each prediction carries the digest of its arm's
  configuration, and a protocol may pin an arm to a digest. If an arm may predict unpinned and be
  pinned later, the digest it is held to can be chosen once its predictions are known.
\item \emph{A reading chosen afterwards.} The comparison becomes readable once the number of
  clusters reaches the pre-registered $N$. A reader who may look at every later count can wait
  for a favourable interval and stop there, the repeated-testing problem of
  \citet{armitage1969}.
\item \emph{An interval that is too narrow.} The default interval is a percentile bootstrap over
  clusters. With few clusters, or clusters that are not independent, it rejects a true null more
  often than its level.
\end{enumerate}

The contributions are:
\begin{itemize}
\item A pin before the first prediction: an arm the primary comparison names is refused a
  prediction until the protocol pins it, no arm with predictions can gain a pin, and a pin is a
  sha256 digest (\cref{sec:pins}).
\item One reading at $N$: once $N$ is fixed the report scores only the first $N$ clusters in the
  order the ledger recorded them, whatever arrives or resolves later, and only once those clusters
  have settled; we state exactly what this does and does not fix (\cref{sec:cut}).
\item Verification that replays the refusal rules, so that a line written outside the append path
  with a valid hash is caught if it breaks one (\cref{sec:pins}).
\item A cluster sign-flip randomization test that a protocol may declare: exact under
  exchangeability of the arms within each cluster with independent clusters, with an add-one
  $p$-value, every sign vector enumerated when $2^k$ fits the protocol's count, and an interval by
  inverting the test. We prove that each sign vector accepts exactly the shifts between the
  weighted means of the clusters it keeps and the clusters it flips, which gives the monotonicity
  the inversion needs and a closed form for its endpoints (\cref{sec:signflip}).
\item Sign vectors built in blocks, so memory no longer grows with the number of vectors
  (\cref{sec:impl}).
\item An evaluation: tests and fourteen mutations, all detected; a simulation study of both
  methods under the null, under dependence across clusters and under an effect; agreement with an
  independent reference on 6{,}000 calls; and memory before and after (\cref{sec:eval}).
\end{itemize}

\section{Background and related work}

\paragraph{The ledger.} Each line holds a sequence number, a kind (protocol, question,
prediction, resolution, void), a body, the time the ledger's own clock recorded it, the previous
line's hash and its own: the sha256 of the canonical JSON of the rest, with 64 zeros before the
first line. Verification recomputes the chain, so editing, dropping or reordering a line breaks it;
removing lines from the end does not, which is why the head is published outside the writer's
reach (a pushed commit, a timestamping service), fixing everything before it. This is the linked time-stamping of \citet{haber1991}. A prediction
carries no time of its own, so an arm cannot backdate one, and a prediction recorded at or after
its question resolved is excluded as late. The scorer reads every number from the protocol: the
blinded pilot reports only the spread of cluster means and the $N$ it implies,
$N = \lceil ((z_{1-\alpha/2} + z_{\text{power}})\, \mathrm{sd} / \delta)^2 \rceil$; the report
refuses to compare before $N$; three controls (a null arm recorded as a copy of a reference, a
derived known edge of size $\delta$ with the run's own noise, and an oracle) tell a real null from
a broken harness or an underpowered run. The decision at $N$ reads one interval: the arm improves
if its lower end is above 0, there is no useful improvement if its upper end is below $\delta$, and
otherwise the result is inconclusive. Scores are Brier's~\cite{brier1950} and the log score, both
strictly proper~\cite{gneiting2007}; secondary comparisons are adjusted by Holm's
procedure~\cite{holm1979}.

\paragraph{Pre-registration and optional stopping.} \citet{nosek2018} argue for fixing the
research question and the analysis before the outcomes are observed, so that a prediction is not a
postdiction presented as one. \citet{armitage1969} showed that testing accumulating data
repeatedly at a fixed nominal level, and stopping at the first rejection, raises the probability of
a false rejection well above that level. The ledger's genesis protocol is the first; the cut at $N$
of \cref{sec:cut} addresses the second by fixing which clusters a reading contains.

\paragraph{Randomization tests.} \citet{fisher1935} and \citet{pitman1937} introduced tests whose
reference distribution is generated by the randomization itself, and \citet{good2005} and
\citet{lehmann2005} treat them in general: when the null distribution of the data is invariant
under a group of transformations, the $p$-value computed over the group's orbit is exact. When the
group is too large to enumerate and elements are drawn at random, counting the observed statistic
among the draws keeps the test exact, which is the add-one rule of \citet{phipson2010}. Inverting a
randomization test for a location shift gives an exact confidence interval; for sign changes of
independent symmetric observations, \citet{hartigan1969} showed that the means of subsamples are
typical values, the intervals between consecutive ordered values each containing the centre with
equal probability. Corollary~\ref{cor:closed} finds the same order statistics, weighted, as the
ends of the inverted interval.

\paragraph{Clustered inference.} The percentile interval of the bootstrap~\cite{efron1993},
resampling whole clusters, is the ledger's default. \citet{cameron2008} document that standard
cluster-robust inference over-rejects with few clusters and propose a wild cluster bootstrap whose
Rademacher weights change the sign of each cluster's residuals, which is a sign flip at the cluster
level. \citet{canay2017} give conditions under which randomization tests over sign changes control
size with a fixed number of clusters when symmetry holds only approximately. Our test is the exact
finite-sample case: it assumes symmetry and independence and does not claim robustness beyond them.

\paragraph{Mutation testing.} A test suite is evidence to the degree that it fails on wrong code;
mutation analysis~\cite{demillo1978} measures this by planting faults. We report fourteen.

\section{Model}\label{sec:model}

No Lean declaration checks the statements below; the propositions on pins and the cut follow from
the ledger's refusal rules, and those on the test are checked against an independent reference in
\cref{sec:eval}.

\subsection{Pins}\label{sec:pins}

\begin{definition}[Protocol in force, pin]
The protocol in force at line $i$ is the last protocol record before $i$. A pin maps an arm to a
digest, a string matching \texttt{(sha256:)?[0-9a-f]\{64\}}. Once any prediction exists, the
primary comparison, $\delta$, $\alpha$, the power, the clip, the resample count, the seed, the
control arms and the method are frozen; $N$ may be fixed once; and no arm, family, secondary
comparison or pin may be dropped.
\end{definition}

The change adds three refusals to the ledger's append path: a protocol whose pins are not all
digests; a prediction from the primary comparison's arm or baseline when the protocol in force does
not pin that arm; and an amendment that adds a pin for an arm that already has a prediction. The
existing rules refuse a prediction whose configuration differs from its arm's pin, and any change
to a pin once predictions exist.

\begin{proposition}[A pin comes first]\label{prop:pins}
In a ledger that verifies, let $a$ be the primary comparison's arm or baseline
and $p$ any prediction of $a$. Then $p$ carries a configuration equal to a pin of $a$ that appears
in a protocol record before the first prediction of $a$, and that pin is the same at every later
line.
\end{proposition}
\begin{proof}[Proof sketch]
The primary comparison is fixed from the first prediction on, so $a$ is named by it at every line
from the first prediction to the end. At $a$'s first prediction the protocol in force pins $a$, or
the prediction is refused, and the configuration must equal that pin. A pin for $a$ cannot be
added after that line, since $a$ then has a prediction, and cannot change, since predictions
exist. So every later prediction of $a$ is held to the same pin.
\end{proof}

The rules are enforced when a record is appended, and verification replays them. It recomputes
the hash chain, then passes each record through the same checks the append path makes, against the
state the earlier records leave and at the time the ledger recorded for that line, never the wall
clock, and fails naming the line and the rule. A line written outside the append path with its
hash recomputed therefore verifies only if it obeys every rule, so \cref{prop:pins} holds of any
file that verifies, not only of files the append path wrote.

\subsection{One reading at N}\label{sec:cut}

\begin{definition}[Pairs, keys and the cut]
A primary pair is a question of the primary family and a round that both primary arms predicted,
whether or not the question has resolved, been voided, or the predictions were late. Its key is the
later of the two predictions' line numbers, and a cluster's key $\kappa(c)$ is the least key of its
pairs. With $N$ fixed, the cut $F_N$ is the set of $N$ clusters of least key among all clusters
with a primary pair. The report reads only the scored rows of clusters in $F_N$, in every table it
prints, and only when every question of the primary family in those clusters is resolved or voided;
otherwise it refuses. A scored row is a prediction whose question is resolved and not voided and
which was recorded before the resolution; its pair has difference
$d = \text{Brier(baseline)} - \text{Brier(arm)}$.
\end{definition}

\begin{proposition}[What waiting cannot change]\label{prop:cut}
Let $L$ be a ledger with at least $N$ clusters having a primary pair, and $L'$ any extension of
$L$ by records the ledger accepts, resolutions and voids included. Then $F_N(L') = F_N(L)$.
\end{proposition}
\begin{proof}[Proof sketch]
A key depends only on the line numbers of two predictions, and a pair's predictions are appended
once. A pair first completed in $L'$ has a key above every line of $L$. So a cluster with a pair in
$L$ keeps its key, the least over a set that gains only larger members, and a cluster first
appearing in $L'$ has a key above the key of every cluster of $L$. The $N$ least keys are
therefore the same clusters. Resolutions and voids change which rows are scored, not any key.
\end{proof}

In particular, neither a cluster that first appears after $N$ is reached nor a late resolution or
void can displace a cluster in $F_N$, however long the reader waits; a voided cluster keeps its
place and contributes no rows. The rows of a cluster in $F_N$ are fixed once the report is allowed
to read: a settled question takes no further prediction, since the ledger refuses one for a
resolved or voided question. One case remains open: a question registered later under the name of
a cluster already in $F_N$ adds a pair, with a key above $K$ and so no change to the cut, but with
rows the cut then reads. \Cref{sec:cutdemo} shows the refusal on a three-cluster ledger.

\subsection{The cluster sign-flip test}\label{sec:signflip}

Let the primary pairs fall into clusters $c = 1, \dots, k$, cluster $c$ holding $n_c \ge 1$
pairs with sum $S_c$ of their differences, and let $n = \sum_c n_c$. The statistic is the
pair-weighted mean $T = \sum_c S_c / n$, the one the bootstrap also reports. The sign vectors
$G = \{-1, +1\}^k$ act on $x \in \mathbb{R}^k$ coordinatewise, and for a shift $\mu$ write
$x_c(\mu) = S_c - \mu n_c$ and
\begin{gather*}
  t_\mu(\sgn) = \Bigl|\sum_c \sgn_c\, x_c(\mu)\Bigr|, \qquad
  p_\mu = \frac{\#\{\sgn \in G : t_\mu(\sgn) \ge t_\mu(\one)\}}{2^k}, \\
  \hat p_\mu = \frac{1 + \#\{r \le R : t_\mu(\sgn^{(r)}) \ge t_\mu(\one)\}}{R + 1},
\end{gather*}
the enumerated and the sampled $p$-values, with $\sgn^{(1)}, \dots, \sgn^{(R)}$ drawn uniformly and
independently from $G$. The denominator $n$ is fixed, so flipping cluster $c$ changes only $S_c$,
and $t_\mu(\sgn)/n$ is the absolute mean of the differences less $\mu$ once the clusters that
$\sgn$ flips are negated. The
test enumerates $G$ when $2^k$ is at most the protocol's count $R$, and samples otherwise.

\begin{definition}[Null hypothesis]
$H_\mu$: the vector $(x_c(\mu))_c$ has the same distribution as $(\sgn_c x_c(\mu))_c$ for every
$\sgn \in G$, given the cluster sizes. It holds at $\mu = 0$ when clusters are independent and the
two arms are exchangeable within each cluster: swapping the arms negates every difference in the
cluster, so each $S_c$ is symmetric about 0.
\end{definition}

\begin{theorem}[Exactness]\label{thm:exact}
Under $H_\mu$, $\Pr(p_\mu \le \alpha) \le \alpha$ and $\Pr(\hat p_\mu \le \alpha) \le \alpha$ for
every $\alpha \in (0, 1)$.
\end{theorem}
\begin{proof}[Proof sketch]
Let $X = x(\mu)$ and let $g_0$ be uniform on $G$ and independent of everything. Under $H_\mu$,
$X$ has the law of $g_0 X$. For enumeration, condition on the orbit $\{gX : g \in G\}$: the
observed statistic is that of a uniformly chosen element, and $p_\mu \le \alpha$ only when it is
among the $\lfloor \alpha 2^k \rfloor$ largest of the $2^k$ orbit values, ties counted against it,
an event of probability at most $\alpha$. For sampling, $(X, \sgn^{(1)}X, \dots, \sgn^{(R)}X)$ has the law of $(g_0X,
\sgn^{(1)}g_0X, \dots)$, and $g_0, \sgn^{(1)}g_0, \dots, \sgn^{(R)}g_0$ are independent and uniform
on $G$, so the $R+1$ statistics are exchangeable given the orbit. The rank of the observed one is
then uniform, ties counted against it, which gives the add-one bound~\cite{phipson2010,lehmann2005}.
\end{proof}

The interval inverts the test: $C_\alpha = \{\mu : p_\mu > \alpha\}$, with $\hat p$ in place of $p$
when the vectors are sampled. Inversion needs the set to be an interval, which the next lemma
provides. For a sign vector $\sgn$ that is not constant, let $P = \{c : \sgn_c = +1\}$ be the
clusters it keeps and $M = \{c : \sgn_c = -1\}$ those it flips, with weighted means
$m_P = \sum_{c \in P} S_c / \sum_{c \in P} n_c$ and $m_M$ likewise.

\begin{lemma}[What one sign vector accepts]\label{lem:accept}
For a non-constant $\sgn$, $t_\mu(\sgn) \ge t_\mu(\one)$ if and only if $\mu$ lies in the closed
interval between $m_P$ and $m_M$. For a constant $\sgn$, $t_\mu(\sgn) = t_\mu(\one)$ for every
$\mu$.
\end{lemma}
\begin{proof}
Write $X_P = \sum_{c \in P} x_c(\mu) = n_P (m_P - \mu)$ and $X_M = n_M (m_M - \mu)$, with $n_P,
n_M > 0$ the pair counts of $P$ and $M$. Then $t_\mu(\one) = |X_P + X_M|$ and $t_\mu(\sgn) = |X_P
- X_M|$, and $(X_P - X_M)^2 \ge (X_P + X_M)^2$ if and only if $X_P X_M \le 0$, that is,
$(m_P - \mu)(m_M - \mu) \le 0$. A constant vector multiplies the sum by $\pm 1$.
\end{proof}

\begin{proposition}[Monotonicity]\label{prop:mono}
With the sign vectors held fixed, $p_\mu$ and $\hat p_\mu$ equal 1 at $\mu = T$ and do not
increase as $\mu$ moves away from $T$ in either direction. Hence $C_\alpha$ is an interval
containing $T$, and bisection on each side finds its ends.
\end{proposition}
\begin{proof}
$T = (n_P m_P + n_M m_M)/n$ lies between $m_P$ and $m_M$, so by Lemma~\ref{lem:accept} the set of $\mu$
each vector accepts is an interval containing $T$. As $\mu$ moves away from $T$ on one side, each
indicator can only switch from 1 to 0, and the $p$-values are increasing functions of their sum.
At $\mu = T$ every indicator is 1.
\end{proof}

The code's docstring argues the same monotonicity by slopes: $t_\mu(\one) = n|T - \mu|$ grows at
rate $n$ away from $T$, while $t_\mu(\sgn)$ changes at rate at most $|\sum_c \sgn_c n_c| \le n$.
Lemma~\ref{lem:accept} gives more: the exact set each vector accepts, and with it the interval in closed
form.

\begin{corollary}[Closed form]\label{cor:closed}
Enumerate $G$, let $m_Q$ range over the weighted means of the $2^k - 2$ non-empty proper subsets
$Q$ of clusters, and let $j = \lfloor \alpha\, 2^{k-1} \rfloor$. If $j = 0$, $C_\alpha =
\mathbb{R}$. Otherwise $C_\alpha = [m_{(j)}, m^{(j)}]$, the $j$-th least and $j$-th greatest of the
$m_Q$. With sampled vectors the ends are the $j'$-th least of $\min(m_P, m_M)$ and the $j'$-th
greatest of $\max(m_P, m_M)$ over the drawn vectors, constant ones counting as unbounded, with
$j' = \lfloor \alpha (R + 1) - 1 \rfloor + 1$.
\end{corollary}
\begin{proof}[Proof sketch]
For $\mu > T$ at most one of $m_P, m_M$ exceeds $T$, so by Lemma~\ref{lem:accept} a non-constant vector
accepts $\mu$ exactly when $m_P \ge \mu$ or $m_M \ge \mu$, and never both. As $\sgn$ ranges over
the non-constant vectors, $P$ and $M$ each range once over the non-empty proper subsets, so $p_\mu =
(2 + 2\,\#\{Q : m_Q \ge \mu\})/2^k$, and $p_\mu > \alpha$ exactly when $\#\{Q : m_Q \ge \mu\} \ge
j$. The lower side is symmetric, and the sampled case counts the drawn vectors directly.
\end{proof}

With unit weights the $m_Q$ are the subsample means that \citet{hartigan1969} calls typical values.
The interval is unbounded exactly when $\alpha < 2/2^k$, the smallest attainable $p$-value.

\begin{proposition}[Coverage under a location shift]\label{prop:cover}
If every difference is $\mu^\ast$ plus noise whose cluster vectors are independent and each
symmetric about 0, then $\Pr(\mu^\ast \in C_\alpha) \ge 1 - \alpha$.
\end{proposition}
\begin{proof}
$H_{\mu^\ast}$ holds, and $\mu^\ast \notin C_\alpha$ exactly when the test at $\mu^\ast$ rejects;
apply Theorem~\ref{thm:exact}.
\end{proof}

The test of $\mu = 0$ needs only $H_0$. The interval needs the location-shift model, in which the
arms' difference moves the whole distribution by a constant; when an effect also changes the shape
of the differences, the interval is approximate while the test of 0 stays exact.

\section{Concepts and intuition}

\paragraph{The digest is the arm.} An arm's predictions are evidence about one configuration. If
the configuration it is held to is chosen after the predictions are known, the experiment can keep
the version that scored well. Requiring the pin first, and forbidding a pin later, makes the
choice before the data; a changed arm becomes a new variant, compared as an arm of its own.

\paragraph{Turn "when" into "which".} A reader who may stop at any count chooses among many
readings. The cut replaces the question of when to read with the question of which clusters to
read, and answers it by the ledger's line order, which no reader controls: the first $N$ clusters,
whatever arrives later.

\paragraph{A flip is a swap.} Under the null, swapping the two arms within a cluster leaves the
cluster's distribution unchanged, and swapping negates the cluster's sum. Every sign vector is
therefore an equally likely version of the data, and the $p$-value asks how extreme the observed
version is among them. Nothing is approximated when every vector is enumerated; when vectors are
drawn, counting the observed one among them keeps the guarantee.

\paragraph{Each vector votes for an interval.} A sign vector splits the clusters into those it
keeps and those it flips. It accepts a shift exactly when the shift lies between the two groups'
weighted means, because only then does flipping move the sum at least as far from zero as the
observed one. The interval is the set of shifts that more than a fraction $\alpha$ of the vectors
accept.

\paragraph{What neither method protects against.} Both methods treat clusters as independent.
Clusters that move together carry less information than their count suggests, so the effective $N$
is smaller and both reject too often. The sign-flip test is exact only where its assumptions hold.

\section{Implementation}\label{sec:impl}

\begin{algorithm}[t]
\caption{The cluster sign-flip test and its inverted interval.}\label{alg:flip}
\begin{algorithmic}[1]
\Require primary pairs with clusters; sign-vector count $R$; seed; level $\alpha$
\Ensure the estimate $T$, the interval $[L, U]$, and the $p$-value at 0
\State $S_c, n_c \gets$ the sum and count of each cluster's differences; $T \gets \sum S / \sum n$
\If{$k < 2$} \Return $T$ with no interval \EndIf
\State $\mathit{exact} \gets 2^k \le R$; \quad $\mathit{rows} \gets 2^k$ if $\mathit{exact}$ else $R$
\State produce the rows in blocks of 1000: row $r$ is the bits of $r$ if $\mathit{exact}$, else
  the next draw from the generator seeded by the seed
\State keep the blocks when $\mathit{rows} \cdot k \le 4 \times 10^6$; otherwise regenerate them
  from the seed at each evaluation
\Function{P}{$\mu$}
  \State $v \gets S - \mu n$; \quad $o \gets |\sum v|$
  \State $h \gets \#\{\text{rows } \sgn : |\sgn \cdot v| \ge o - 10^{-12} \sum |v|\}$, block by block
    \label{line:tol}
  \State \Return $h / \mathit{rows}$ if $\mathit{exact}$ else $(1 + h)/(\mathit{rows} + 1)$
\EndFunction
\Function{Edge}{$s$} \Comment{$s = \pm 1$}
  \State $\mathit{step} \gets \max(10^{-3}, \max_c |S_c / n_c|)$; \quad $\mathit{in} \gets T$;
    \quad $\mathit{out} \gets T + s \cdot \mathit{step}$
  \State double $\mathit{step}$ until $\Call{P}{\mathit{out}} \le \alpha$, at most 60 times,
    else \Return $s \cdot \infty$
  \State bisect 50 times, keeping $\Call{P}{\mathit{in}} > \alpha \ge \Call{P}{\mathit{out}}$
  \State \Return $\mathit{in}$
\EndFunction
\State \Return $T$, $[\Call{Edge}{-1}, \Call{Edge}{+1}]$, $\Call{P}{0}$
\end{algorithmic}
\end{algorithm}

\Cref{alg:flip} is the test as written. The tolerance in line~\ref{line:tol} counts as ties the comparisons that
floating-point rounding breaks, chiefly the two constant vectors, which tie exactly in real
arithmetic. Blocks are produced in a fixed order from one generator, so the drawn vectors, and
every result, are the same whatever the block size; a test checks this with blocks of 100 rows and
regeneration forced. Before the change, all vectors were built at once as a $\mathit{rows} \times
k$ array of 64-bit floats, with integer intermediates of the same shape, and the protocol's count
has no upper bound.

The cut is computed from the table the scorer builds from the chain: each primary pair's key is the
later line of its two predictions, whatever the state of its question, each cluster's key the
least of its pairs', and the table carries the clusters in key order and those with a question not
yet resolved or voided. When more than $N$ clusters have pairs the report keeps the rows of the $N$
of least key, refuses while one of them is unsettled, and reports how many clusters and pairs it
left out. Verification is the replay of \cref{sec:pins}, and the \texttt{score} and \texttt{power}
commands run it first. The pin rules are three checks in the append path, described in \cref{sec:pins}. The
method is a protocol field, \texttt{percentile} by default so that existing ledgers score as
before, and frozen with the rest of the bar once a prediction exists.

\section{Evaluation}\label{sec:eval}

\subsection{Setup}
All runs used macOS 27.0.1 on an Apple M2 Pro with 32\,GiB of memory, Python 3.15.2 and NumPy
2.5.4, on an export (\texttt{git archive}) of the branch at \texttt{0d7cb89}, the simulations and memory
figures below; the pre-change memory figures use the commit before the blocks change,
\texttt{caf5c0a}. The tests and mutations were run on the later head \texttt{62d8187}, which changes
verification and the cut and none of the simulated procedures. Simulations call the
branch's own \texttt{cluster\_signflip} and \texttt{cluster\_bootstrap} with the protocol default
of 10{,}000 sign vectors or resamples, eight worker processes, and one seeded generator per
replication, so every number below can be reproduced from its description.

\subsection{Tests and mutations}
The two test files hold 80 tests (43 for the ledger, 37 for the scorer), which pass in 4\,s.
\Cref{tab:mutations} lists fourteen single edits and the tests each fails, and all fourteen are
detected. Two of them, dropping the add-one rule and flipping pairs instead of clusters, survived
an earlier version of the suite: the first changes the $p$-value by at most 0.1 percentage points at
1{,}000 vectors, too little for a rate test to see, and the second cannot be told from the right
test on generated differences with no within-cluster correlation. They are now killed by direct
cases: a hand-built block of sign vectors where the add-one $p$-value is $3/5$ and the plain one
$2/4$, and three clusters of four identical pairs, where flipping clusters gives $p = 2/8$ and
flipping pairs $2/4096$. The cut and verification mutations were added with their fixes; of the
earlier cut mutations only the counts of failing tests changed. Passing is a statement about these
edits and no more.

\begin{table}[ht]
\centering\small
\caption{Single mutations of the ledger and scorer, each run against the 80 tests.}
\label{tab:mutations}
\begin{tabular}{@{}lrl@{}}
\toprule
Mutation & Fail & Failing tests \\
\midrule
a primary arm may predict unpinned & 1 & the pin-first test \\
a pin may be added after the arm predicted & 1 & the no-retrofit test \\
any string is a pin & 7 & the seven malformed pins \\
the method is not frozen & 1 & the frozen-method test \\
no cut at $N$ & 5 & the five tests of the cut \\
the cut by cluster name, not ledger order & 2 & waiting past $N$; first $N$ by ledger order \\
no add-one when vectors are sampled & 1 & the add-one case \\
always sample, never enumerate & 1 & the three-cluster exact $p$-value \\
a one-sided statistic & 1 & the same \\
blocks reseeded, not one stream & 1 & the block-size test \\
pairs flipped instead of clusters & 1 & the per-cluster case \\
verification checks the chain only & 2 & the forged unpinned and forged late lines \\
the cut among resolved clusters only & 2 & late resolution; void keeps its place \\
no wait for open clusters in the cut & 1 & late resolution \\
\bottomrule
\end{tabular}
\end{table}

\subsection{The methods under the null}
Data were generated with two exchangeable arms: each cluster has one event of probability drawn
uniformly from $[0.1, 0.9]$ and one to three questions turning on it; each arm predicts the event
probability plus an arm-by-cluster error of standard deviation 0.10 and an arm-by-question error of
0.05, clipped to $[0.01, 0.99]$; $d$ is the Brier difference. The null holds exactly, so every
rejection is false. \Cref{tab:null} gives the share of 10{,}000 replications in which each method's
interval excluded 0 (for the sign-flip test, the same as $p \le \alpha$, which a test checks). With
13 clusters the sign-flip test enumerates all $2^{13}$ vectors. The percentile bootstrap rejects
about twice as often as its level there, and 1.3 to 1.9 points too often at 40 clusters; the
sign-flip test stays within Monte Carlo error of its level (one standard error is 0.22 points at
5\% and 0.16 at 2.5\%).

\begin{table}[ht]
\centering\small
\caption{False rejections under the null, 10{,}000 replications each.}\label{tab:null}
\begin{tabular}{@{}lrrrr@{}}
\toprule
& \multicolumn{2}{c}{$\alpha = 0.05$} & \multicolumn{2}{c}{$\alpha = 0.025$} \\
\cmidrule(lr){2-3}\cmidrule(lr){4-5}
Clusters & sign-flip & bootstrap & sign-flip & bootstrap \\
\midrule
13 (enumerated) & 4.85\% & 9.67\% & 2.62\% & 6.29\% \\
40 (sampled) & 5.05\% & 6.86\% & 2.32\% & 3.84\% \\
\bottomrule
\end{tabular}
\end{table}

For coverage the data follow the location-shift model of Proposition~\ref{prop:cover}: $d = 0.1 + u_c +
e_{cj}$ with $u_c \sim \mathcal{N}(0, 0.1^2)$, $e_{cj} \sim \mathcal{N}(0, 0.2^2)$, one to three
pairs per cluster and 40 clusters. Over 10{,}000 replications the sign-flip interval covered 0.1 in
94.90\% of them at 95\% and 97.38\% at 97.5\%; the bootstrap's covered 93.79\% and 96.31\%.

\subsection{Dependence across clusters}
To model residual dependence, clusters were grouped in blocks of four sharing an effect $w \sim
\mathcal{N}(0, \tau^2)$, added to every difference in the block, with $u_c$ and $e_{cj}$ as above
and no true effect. For $\tau = 0.05$ and $0.10$ the correlation between the means of two clusters
in a block is about 0.07 and 0.22. \Cref{tab:dep} shows that both methods over-reject and that the
excess grows with $\tau$; the sign-flip test's excess is the smaller at 40 clusters, and the two
nearly agree at 120.

\begin{table}[ht]
\centering\small
\caption{False rejections with clusters dependent in blocks of four, 4{,}000 replications
each.}\label{tab:dep}
\begin{tabular}{@{}llrrrr@{}}
\toprule
& & \multicolumn{2}{c}{$\alpha = 0.05$} & \multicolumn{2}{c}{$\alpha = 0.025$} \\
\cmidrule(lr){3-4}\cmidrule(lr){5-6}
Clusters & $\tau$ & sign-flip & bootstrap & sign-flip & bootstrap \\
\midrule
40 & 0 & 5.28\% & 6.65\% & 2.78\% & 3.75\% \\
40 & 0.05 & 7.68\% & 9.45\% & 3.68\% & 5.30\% \\
40 & 0.10 & 14.08\% & 16.15\% & 9.18\% & 11.35\% \\
120 & 0.05 & 7.93\% & 8.35\% & 4.45\% & 4.98\% \\
120 & 0.10 & 12.90\% & 13.50\% & 8.53\% & 9.18\% \\
\bottomrule
\end{tabular}
\end{table}

\subsection{Power at the planned N}
The pilot's formula was applied to the location-shift model with $\delta = 0.05$. The standard
deviation of a cluster's mean difference, estimated from 200{,}000 simulated clusters, is 0.186,
which gives $N = 109$ clusters for 80\% power at $\alpha = 0.05$ and $N = 132$ at $\alpha = 0.025$.
\Cref{tab:power} gives the share of 4{,}000 replications in which the interval's lower end exceeded
0, the verdict that the arm improves, at the true shift. With $N$ derived for the level in use,
both methods reach the planned 80\%; keeping $N = 109$ at $\alpha = 0.025$ costs about nine points.
The bootstrap's extra point of power goes with its excess false rejections (\cref{tab:null}). With
the block dependence of $\tau = 0.05$ at $N = 132$, power fell to 77\%, and the sign-flip interval
missed the true shift in 4.0\% of replications against a nominal 2.5\%.

\begin{table}[ht]
\centering\small
\caption{Power at the true shift $\delta = 0.05$, 4{,}000 replications each.}\label{tab:power}
\begin{tabular}{@{}rlrrr@{}}
\toprule
Clusters & Dependence & $\alpha$ & sign-flip & bootstrap \\
\midrule
109 & none & 0.05 & 81.7\% & 82.9\% \\
109 & none & 0.025 & 72.9\% & 74.6\% \\
132 & none & 0.025 & 82.3\% & 83.3\% \\
132 & $\tau = 0.05$ & 0.025 & 77.3\% & 78.2\% \\
\bottomrule
\end{tabular}
\end{table}

\subsection{Agreement with an independent reference}
The reference implements Theorem~\ref{thm:exact} and Corollary~\ref{cor:closed} directly. In the enumerated branch it counts
all $2^k$ vectors in exact rational arithmetic and takes the interval's ends as order statistics of
the subset means; in the sampled branch it draws the same vectors from the seed in one call and
takes the ends from Corollary~\ref{cor:closed}. Over 3{,}000 data sets per branch (2 to 13 clusters
enumerated, 14 to 60 sampled with 1{,}000, 2{,}500 or 10{,}000 vectors, $\alpha$ of 0.10, 0.05 or
0.025, one to four pairs per cluster), the branch's function and the reference gave the same
$p$-value on all 6{,}000 calls and the same ends to within a relative $7.2 \times 10^{-12}$. A first
version of the reference compared floating-point sums without a tolerance and disagreed on 19
sampled $p$-values, each by the drawing of a constant vector whose floating sum rounded below the
observed one; counted as the exact ties they are, as the code's tolerance counts them, the
disagreements vanished. The bisection therefore finds the closed-form ends, and the closed form
itself (Corollary~\ref{cor:closed}) could replace it in the enumerated case.

\subsection{Memory and time}
\Cref{tab:memory} gives the peak resident memory of a process making one call, and the call's wall
time, before and after the blocks change, on synthetic differences with one to three pairs per
cluster. The $p$-value and both ends of the interval are identical before and after in every case.
After the change the peak is the interpreter's own footprint at every size. The price is time:
above $4 \times 10^6$ sign entries the vectors are regenerated from the seed at each of the hundred
or so evaluations, between 4 and 14 times slower than holding them. At the default of 10{,}000
vectors and up to 400 clusters they are kept, and the call is unchanged.

\begin{table}[ht]
\centering\small
\caption{One call of the test: peak memory of the process and wall time.}\label{tab:memory}
\begin{tabular}{@{}lrrrrr@{}}
\toprule
& & \multicolumn{2}{c}{before (\texttt{caf5c0a})} & \multicolumn{2}{c}{after (\texttt{0d7cb89})}
\\
\cmidrule(lr){3-4}\cmidrule(lr){5-6}
Sign vectors & Clusters & memory & time & memory & time \\
\midrule
$2^{20}$, enumerated & 20 & 560\,MB & 0.85\,s & 39\,MB & 4.3\,s \\
$2^{23}$, enumerated & 23 & 4.81\,GB & 9.5\,s & 40\,MB & 39.4\,s \\
$10^5$, sampled & 200 & 362\,MB & 0.50\,s & 51\,MB & 6.5\,s \\
$10^6$, sampled & 200 & 3.26\,GB & 4.8\,s & 51\,MB & 64.8\,s \\
$10^4$, sampled & 132 & 59\,MB & 0.02\,s & 50\,MB & 0.02\,s \\
\bottomrule
\end{tabular}
\end{table}

\subsection{The cut, demonstrated}\label{sec:cutdemo}
On a ledger with $N = 2$ and three clusters, recorded in the order $x, y, z$, with $y$ and $z$
resolved before $x$, the earlier cut, made among resolved clusters, read $\{y, z\}$ with mean
difference 0.70 once $y$ and $z$ resolved, and $\{x, y\}$ with mean 0.00 after $x$ resolved,
reporting $z$ as beyond $N$: the reading moved after $N$ was reached. The cut now ranges over all
three clusters whatever their state. With $y$ and $z$ resolved and $x$ open the report refuses,
naming the open cluster among the first two; when $x$ resolves it reads $\{x, y\}$ and reports one
cluster beyond $N$. Resolving $z$ late, with $x$ and $y$ already read, leaves the primary and
per-arm tables unchanged, and voiding $x$ leaves it in its place, so that only $y$ is scored and
$z$ stays beyond $N$. The test for waiting adds twenty clusters after the thirtieth with all
questions resolved, and the report is unchanged.

\subsection{The audit that motivated the change}
An audit of an experiment built on the ledger reported the figures that led to the method; its data
and dependence model belong to the audited application and are not reproduced here. On pure
symmetric noise the implementation rejected 4.8\% at $\alpha = 0.05$ against 6.5\% for the
percentile bootstrap on the same data, and 2.7\% at $\alpha = 0.025$; its interval covered 94.75\%;
and it matched the audit's own reference implementation in all 36{,}000 calls. Under the audit's
model of realistic dependence the percentile bootstrap rejected 5.9\% with independent clusters and
7.5\% to 11.4\% with residual dependence, while the sign-flip test at $\alpha = 0.025$ held 2.9\% to
5.9\% at the $N$ re-derived for that level, with 82\% power. \Cref{tab:null,tab:dep} reproduce the
direction and rough size of these figures on generators stated in full.

\section{Limitations}\label{sec:limits}

\paragraph{Replay does not authenticate.} Verification checks that every line obeys the rules, not
that the writer is who they say. A line that obeys the rules and is added after the last published
head, by anyone who can write the file, verifies. Against such a writer the guarantee is the
anchored head: lines before a published head cannot be changed. Replay also reads every older
ledger against today's rules, so a ledger written before the digest rule, with a shorter pin,
would now fail to verify.

\paragraph{A pin is a digest of something.} The ledger checks that a pin has the form of a sha256
digest, not that it is the digest of the configuration the arm runs; that binding is the arm's.

\paragraph{A cluster can still gain a question.} The cut is fixed and the report waits for the
questions of its clusters to settle, but a question registered later under the name of a cluster in
the cut adds rows to it (\cref{sec:cut}). Refusing a new question in a cluster that is already in
the cut would close this; it is not done.

\paragraph{Independence across clusters.} Neither method protects against clusters that move
together; under the block model the sign-flip test rejected 7.7\% to 14.1\% at $\alpha = 0.05$
(\cref{tab:dep}). Its exactness is a property of the assumptions, which the experiment must make
true by how it forms clusters.

\paragraph{The interval assumes a location shift.} The test of 0 is exact under $H_0$ alone; the
interval's coverage is exact only when the effect shifts the differences by a constant.

\paragraph{Resolution and cost.} With $R$ sampled vectors the smallest $p$-value is $1/(R+1)$, and
the interval takes about a hundred evaluations of the test, each over all vectors; above the
keeping threshold each evaluation regenerates them, which bounds memory at a cost of 4 to 14 times
the time (\cref{tab:memory}). The planning formula for $N$ is a normal approximation for the mean of
cluster means and does not depend on the method.

\paragraph{Mutation testing is a sample.} The fourteen edits are the ones we thought of; a suite
that kills them can still miss another.

\section{Future work}\label{sec:future}
Refusing a question that joins a cluster already in the cut would make the rows, as well as the
membership, of the reading fixed. When the dependence between clusters has a known structure,
flipping blocks of clusters together restores exactness at the block level, at the cost of fewer
units. In the enumerated case, Corollary~\ref{cor:closed} computes the interval from one sort of
the subset means instead of about a hundred passes over the vectors.
Lemma~\ref{lem:accept} and Corollary~\ref{cor:closed} are elementary enough to state and check in
Lean.

\section{Conclusion}
A pre-registered bar can still be bent after the data are seen by choosing an arm's configuration,
the moment of reading, or an interval that is too narrow. The ledger now requires a pin before an
arm's first prediction and forbids adding one later, verifies by replaying its refusal rules, reads
only the first $N$ clusters in its own line order, whatever resolves later, and offers a sign-flip test that is exact when clusters are independent and the arms
exchangeable within them. Each sign vector accepts the shifts between the weighted means of the
clusters it keeps and flips, which makes the inverted interval well defined and gives its ends in
closed form, and an independent reference built on that form agrees with the implementation on
every call. The evaluation also finds what the change leaves open: a cluster that can still gain a
question after the cut, a verification that checks obedience and not authorship, and no method that
tolerates dependence across clusters.

\bibliographystyle{plainnat}
\bibliography{refs}
\end{document}
